\section{功能模块与业务能力}
\subsection{六类校园服务}
当前运行组合将五类 MCP 校园服务与校车本地 Provider 接入统一能力层。这里的“六类”按面向用户的业务领域统计，不等于 MCP 工具数量，也不包括控制台中仅登记的可选服务。

\subsubsection{评课社区}
评课能力支持课程与教师搜索、课程详情、评论、排名及统计。高层查询入口 \code{icourse.public-query.v2} 将业务操作与查询条件映射至 iCourse 服务。回答需要区分同名课程、不同教师、评论条数与站点统计；评论中的主观体验应表述为评论观点，不能转化为对课程或教师的确定性评价。

在选课场景中，多位群友可以分别提出讲解细致程度、证明训练和作业负担等偏好。工具提供课程与评论材料，回答模型负责按这些维度组织建议。若评论没有覆盖某一维度，应说明材料不足；学期是否实际开课属于教务信息，不能仅靠评课社区推断。

\subsubsection{第二课堂}
二课服务提供活动搜索、详情、类别字典及连接状态。参数绑定处理活动关键词、类别、时间和数量，返回活动名称、时间、报名状态与资料范围。运行环境通过部署端 CAS 凭据访问服务；群聊问答面向活动信息，个人学时、已报名记录、签到及代报名不属于本次智能体查询范围。

自然语言中“最近”“公开发布”“学术活动”等条件需要分别落实为查询或结果过滤。修复记录表明，类别未绑定、活动详情参数错误或返回院系树过大，都会导致错误内容或过高的上下文消耗。项目针对这些具体问题收敛了参数和返回结构。

\subsubsection{教务公开信息}
教务服务支持学期列表、全校开课、考试和教学日历。课程名称与教师来自开课记录，学期和日期字段随结果保留。检索失败、空结果与没有考试安排是不同情况，回答应按实际返回区分。当前接口不提供学生个人成绩查询，也不替代校方教务操作。

\subsubsection{培养方案研究资料}
培养方案服务按年级、专业、课程代码和方案类型查询，并支持方案与课程比较。数据来自公开资料整理的非官方研究快照，当前来源覆盖 2015—2026 级，快照日期为 2026 年 8 月 2 日。课程身份使用精确课程代码，防止同名课程在不同方案中被误合并。

该服务适合查阅和讨论培养要求。报告中的学分结果必须绑定年级、专业与方案类型；不能把快照当成实时教务数据库，或把一次查询当成个人毕业资格审核。

\subsubsection{校园通知}
NotifAI 服务提供通知搜索、详情、日历、截止日期、来源与分类字典及统计。不同操作拥有独立能力定义，查询内容与查询字典应分开。回答保留标题、发布日期、截止时间和来源链接，在时间字段缺失时明确说明。

“最近三条通知”属于通知正文检索，不能误选“列出通知来源”。相关录制回归针对这类口语到工具操作的误选作了修复，说明工具成功返回并不自动意味着业务选择正确。

\subsubsection{校车查询}
校车采用本地版本化时刻表，通过 \code{ustc.shuttle.public-query.v1} 进入 Runtime。插件不注册独立自然语言消息处理器，意图识别、查询、总结和发送仍由核心负责。查询支持校区、路线、日期类型和时段，返回相应班次。

当前资料中东区、西区、南区及高新校区版本日期为 2026 年 8 月 30 日，太湖路园区为 8 月 27 日。实现区分始发站与经过站时间，并保留班次备注；“没有公交车辆运行”不等于该班次取消。时刻表通过维护数据文件更新，属于静态资料查询。

\subsection{近期对话、协作与表达}
直接交流可以用于解释问题、分析方案、整理多人分工和回答上下文追问。近期讨论总结从可获取的同群历史构建，控制台显示读取条数和覆盖范围；若只有最近若干条消息，就不应称为“全天完整摘要”。分工草案需要将人员、任务和产出对应起来，不能只列抽象角色。

Persona 资产使同一智能体在技术讨论与日常交流中保持连贯语气。精确输出要求优先于人格表达：要求只回复指定文字时不得附加语气词，要求 JSON 时不得插入解释或表情。长回答由输出层按平台能力处理。

\subsection{群插件与可选服务}
\begin{longtable}{@{}L{31mm}L{51mm}L{67mm}@{}}
\caption{扩展资产与当前使用范围}\\\toprule
组件 & 已有功能 & 本次状态与限制\\\midrule\endfirsthead
\toprule 组件 & 已有功能 & 本次状态与限制\\\midrule\endhead
Emoji Kitchen & 显式表情组合，返回合成图像 & 已有真实服务生成素材；与语言模型生成图片分开。\\
复读插件 & 根据重复消息规则响应 & 独立、按群配置；录制固定样例验证处理器，不代表长期自然群表现。\\
Arcaea 扩展 & 曲目信息、本地谱面、绑定及 B50 相关资产 & 曲目与谱面有已生成素材；真实 B50 查分本次保持关闭；本地 B50 Provider 未接入生产组合。\\
Sub2API Readonly & 显式服务统计命令 & 录制使用真实处理器处理固定统计样例，不能将其当成实时业务统计。\\
社交扩展插件 & 生日、睡眠记录、投票等显式命令与状态组件 & 默认关闭，独立于 Core；不自动发送生日祝福或情绪安慰。\\
五个可选 MCP & 本地推荐、专业设置、学校通知、学院通知、图书馆开放时间 & 已有服务与目录接入、连接检测记录；未配置相应 Capability mapping 时，不进入当前 Planner 候选。\\
旧兼容组件 & Target Talk、Reply Polish、旧记忆等源码与迁移材料 & 不据源码存在推断当前启用；当前主链由 Core 负责入站与输出。\\\bottomrule
\end{longtable}

将独立插件、可发现服务和可规划能力分开描述，可以明确回答“功能已写好”“服务能连接”“用户可以通过自然语言调用”这三个不同问题。扩展资产为后续接入提供基础，本次成果以已经贯通的业务路径为主。
\source{[S5, S6, S7, S10]；各校园服务与插件 README，原子查询和素材记录。}
